\begin{mse}
    \escheme{Проведено измерение ЭДС батареи}{0.5}

    \escheme{Найден ток при открытом диоде}{0.7}

    \escheme{Найдено сопротивление первого резистора}{0.5}

    \escheme{Найден ток при закрытом диоде, то есть при другой полярности}{0.7}

    \escheme{Утверждение и аргументация последовательного соединения диода и резистора к концам черного ящика}{1.0}

    \escheme{Утверждение и аргументация о том, что к выводам черного ящика могут быть соединены лишь диоды с резистором или ветви с конденсатором, обеспецивающие разрыв цепи}{0.6}

    \escheme{Описано поведение/наблюдение подтверждающее присутствие конденсатора}{0.5}

    \escheme{Описан метод измерения напряжения на выводе при разрядке конденсатора или нарисована соответствующая схема, разъясняющая ход эксперимента}{0.8}

    \blockbegin{Построена таблица зависимости $V(t)$. Баллы за измерения ставятся только если сами измерения и метод верные.}{1.8}{7}
        \bscheme{Измерено 15 и больше точек напряжения}{0.9}
        \bremark{Между 10 и 14 точками}{0.6/0.9}
        \bremark{Между 5 и 9}{0.3/0.9}
        \bscheme{Есть измерения при $t > \qty{20}{\s}$}{0.3}
        \bscheme{Есть измерения при $t < \qty{6}{\s}$}{0.3}
        \bscheme{Вычислены значения $\ln V$ для полученных точек}{0.3}
    \blockend

    \blockbegin{Построен график зависимости $V(t)$. Баллы за пункты \textbf{M10} и \textbf{M11} ставятся только если принят пункт \textbf{M9}.}{1.8}{4}
        \bscheme{Оси подписаны и оцифрованы}{0.5}
        \bscheme{Нанесены все точки в соответствии с таблицей}{0.9}
        \bscheme{Проведена сглаживающая кривая}{0.4}
    \blockend

    \blockbegin{Построен график в координатах $\ln V, t$}{2.3}{4}
        \bscheme{Оси подписаны и оцифрованы}{0.5}
        \bscheme{Нанесены все точки в соответствии с таблицей}{0.8}
        \bscheme{Кривая аппроксимирован прямой функцией}{0.5}
        \bscheme{Из графика найдены значения $k, b$}{2$\times$ 0.3}
    \blockend

    \escheme{Предложена и обоснована схема черного ящика}{1.3}

    \escheme{Формула \eqref{2}:  $\ln V = \ln V_0 \frac{R_1}{R_1+R_2} - \frac t \tau$}{0.8}

    \escheme{Найдено сопротивление второго резистора}{0.4}

    \escheme{Найдена емкость конденсатора}{0.9}

    \etotal{15.0}
\end{mse}
